\documentclass{article}

\usepackage{multicol}
\usepackage{times}
\usepackage{graphicx}

\textheight 9in
\textwidth 6.5in
\oddsidemargin 0in
\evensidemargin 0in
\topmargin -.5in
\pagestyle{empty}
\newcommand{\setspace}[1]{\renewcommand{\baselinestretch}{#1}\large\normalsize}
\setspace{1.05}
\renewcommand{\textfraction}{0.1}
\renewcommand{\topfraction}{0.9}
\newcommand{\tods}{{\em TODS\/}}

\begin{document}
\begin{center}
{\Large\bf Rights of {\em TODS} Readers, Authors and Reviewers}\\[2ex]
{\large\it  Richard Snodgrass}\\[1ex]
{\tt rts@cs.arizona.edu}\\[3ex]
\end{center}

About eighteen months ago the ACM Publications Board approved a sweeping
policy, termed the ``Rights and Responsibilities in ACM Publishing'' ({\tt
http://www.acm.org/pubs/rights.html}), summarizing the rights and
responsibilities of readers, authors, reviewers, editors, program chairs and
committees, and libraries vis-a-vis journals, transactions, magazines,
conference proceedings, and SIG newsletters published by ACM~\cite{SN02}.

This policy document contains a mixture of current practice and goals for the near
future. I and the \tods\ Editorial Board take these stated rights very
seriously and have been working for over a year on developing mechanisms,
policies, and procedures to ensure each of these rights. As a result of
these efforts, \tods\ is the \underline{first} ACM transactions to fully
implement all 42 (!) rights listed here.

I'll discuss the steps we have and are taking. (I'll not get into
rights of Editorial Board members, except to state that those rights have
also been fully implemented.)

Included are several new policies that you as a prospective author or reviewer may
appreciate.
\begin{itemize}
\item \tods\ now publishes its turn-around statistics.
\item \tods\ now limits reviewers to one review per year.
\item \tods\ now gives each reviewer at least two months to perform a review.
\end{itemize}

\section{Readers}

\noindent
Quotations from the ACM Rights and Responsibilities document are given in {\em italics}.

\vspace{2ex}

\noindent
{\em Readers can expect ACM to}

\begin{itemize}
\item {\em Publish on time with the printed and Digital Library versions
  available the first day of the issue month}

  Readers know that for several years each issue of \tods\ has come out
  months after its cover date. As I discussed in my September 2001
  editorial ({\tt http://doi.acm.org/10.1145/502030.505049}),
  there were two primary underlying causes: an inadequate backlog 
  and slow production. I'm pleased to report that both problems
  have been solved, and that the December 2002 issue will be available in print
  and in the digital library on December 1, 2002. In fact, articles are
  placed on the \tods\ web site as soon as they are accepted. The December
  articles have been present there since September.

\item {\em Ensure that articles are accurate and of high quality}

  The \tods\ Editorial Board does its part through rigorous reviewing of
  submitted manuscripts, and ACM publication staff do their part through
  careful copy-editing and close interaction with authors (see also
  {\tt http://www.}\linebreak{\tt acm.org/pubs/quality.html}).

\item {\em Ensure that the electronic and printed version of an article
  match within the limits of the style guidelines of each format}

  Starting a few issues ago, the electronic and printed versions of each
  \tods\ article were and continue to be
  generated from the same \LaTeX\ source, to ensure that the contents match.

\item {\em Ensure that journal, transactions, and magazine articles are
  professionally copy-edited}

  This is done for every \tods\ article.

\item {\em Ensure consistent formatting of articles in each publication}

  Authors of accepted manuscripts are provided with style files (see {\tt
  http://www.acm.org/pubs/sub}\-{\tt missions/submission.htm}) to ensure consistent
  formatting. The final version the authors provide is identical to that
  which appears, down even to the fonts employed, modulo subsequent copy editing.

\item {\em Make publications available at low cost to individual
  subscribers, for the current year and for all previous years}

  \tods\ is available to individual subscribers for \$39 per year (four
  issues) for ACM Members, \$34 per year for ACM Student Members, and \$164
  for nonmembers. An electronic subscription is \$31 for members, \$27 for
  student members, and \$131 for non-members. A combined print and online
  subscription is \$47 for ACM Members, \$41 for ACM Student Members, and \$197
  for nonmembers. The online subscription includes all previous years: the
  ACM Digital Library contains all issues back to issue 1 of volume 1, published in
  March 1976.

\item {\em Take into account the needs of readers in economically emerging
  countries and in economically undeveloped countries}

  Like all ACM publications, \tods\ has a sliding subscription rate for
  such readers.

\item {\em Enable fast access to the electronic version of each article,
  throughout the world}

  ACM has invested considerable funds so that all readers have fast access.

\item {\em Permit low cost  purchasing of individual copies of articles
  (printed or electronic version).}

  Single print copies are \$18 each to ACM members and \$40 to nonmembers;
  single electronic copies are \$5 and \$10, respectively.
\end{itemize}

\section{Authors}

The policy partitions the right of authors into three parts, corresponding
to the phases a submitted paper goes through: reviewing, processing of
accepted works, and dissemination.

\subsection{Reviewing}

\noindent
{\em When an author makes a submission, a confidential review process is
initiated. The aim of the review process is to make an appropriate and
timely decision on whether a submission should be published. ... Thus
authors can expect ACM to}

\begin{itemize}
\item {\em Keep them informed on the status of their submission}

  Currently communication of progress of the review of a submission is
  handled via email from the Associate Editor handling the submission.

  \tods\ is implementing an online manuscript tracking system that will
  afford authors more information on the status of the processing of
  manuscripts they have submitted. 

\item {\em Use impartial reviewers}

  This has been a policy of \tods\ from the beginning.

\item {\em Issue timely review and clear feedback}

  Reducing the turn-around time for submitted manuscripts has been a priority.
  For all manuscripts
  submitted between July 1, 2001 and March 30, 2002 (the latest date for
  which statistics are available, as I write this on October 2, 2002), the
  average time required to process a manuscript, from submission to
  editorial decision, is 4.3 months. This interval
  has been quite constant over the past six months or so.

  Also of interest is the {\em maximum} turn-around time. This value had jumped
  around a lot, but is now settling down. We are committed to bounding the
  maximum turn-around time to 6 months.

\item {\em Maintain confidentiality.}

  This has also been a policy of \tods\ from the beginning.

\end{itemize}

\subsection{Processing of Accepted Works}

\noindent
{\em Once a submission has been accepted, authors can expect ACM to publish the
work in a timely and professional manner. Authors can expect to have
approval of all changes to the work. In addition, ACM will strive to not
cause authors to perform unnecessary work .... Thus authors can expect ACM to}

\begin{itemize}
\item {\em Provide reasonable time to fix galleys for journal, transactions,
  and magazine articles}

  \tods\ keeps authors abreast of where their paper is in the production
  process, and gives authors at least a week to check proofs.

\item {\em Note all copy-editing changes for journals and transactions}

  The production process for \tods\ provides a marked-up copy of the
  original draft, along with the updated copy, so that authors know exactly
  what changes were made to their paper.

\item {\em Seek author approval of the final copy for journal, transactions,
  and newsletter articles}

  This has been a policy of \tods\ at least with the new production process
  instituted several months ago.

\item {\em Not introduce errors in the production process}

  Now that \tods\ uses \LaTeX\ (the formatter of choice for most
  submissions) internally in its production process, the number of errors
  introduced has gone down significantly, as most papers need not be
  converted.

\item {\em Add no material without the corresponding author's approval}

  ACM always checks with authors on copy-editing changes to \tods\ articles.

\item {\em Be financially responsible for its own internal preparation
  costs}

  The subscription rate for \tods\ has been kept low by careful selection of
  vendors by the ACM publications staff.

\item {\em Ensure metadata accuracy for journal, transactions, and magazine
  articles.}

  ACM goes to great lengths to ensure that the metadata in its Digital
  Library is accurate.
\end{itemize}

\subsection{Dissemination}

\noindent
{\em Publication is only a part of the broader goal of disseminating ideas and
results.  Authors can expect ACM to contribute to this wider goal, and in
particular to encourage dissemination in multiple forums. ... Thus authors can
expect ACM to} 

\begin{itemize}
\item {\em Allow a submission to be posted on home pages and public
  repositories before and after review}

  The ACM Copyright Policy ({\tt
  http://info.acm.org/pubs/copyright\_policy/}) is quite explicit (and
  author-friendly) on this point: ``As part of their retained rights,
  authors may revise their ACM-copyrighted work and post the new version on
  non-ACM servers (personal, employer, or CoRR.)''

\item {\em Allow an authors' version of their own ACM-copyrighted work on
  their personal server or on servers belonging to their employers}

  The copyright policy is also explicit on this point. ``Authors may post an
  author-version of their own ACM-copyrighted work on a personal server or
  on a server belonging to their employer, but they may not post a copy of
  the definitive version that they downloaded from ACM's Digital
  Library.''

\item {\em Allow metadata information, e.g., bibliographic,
  abstract, and keywords, for their individual work to be openly available}

  This holds for all ACM publications.

\item {\em Allow authors the right to reuse their figures in their own
  subsequent publications for which they have granted ACM copyright}

  This holds for all ACM publications.

\item {\em Provide statistics for each journal, transaction, and newsletter
  on its average turn-around time and its current backlog of articles.}

  \tods\ is the first database journal, the first ACM journal, and indeed
  the first journal that I am aware of, that publishes its turn-around
  performance (see {\tt http://www.acm.org/tods/TurnaroundTime.html}) for
  all to see.

  The current backlog of articles is also maintained (at
  {\tt http://www.acm.org/tods/Upcoming.html}). Authors can judge for
  themselves how responsive the journal is.
\end{itemize}

\section{Reviewers}

\noindent
{\em ACM recognizes that the quality of a refereed publication rests primarily on
the impartial judgment of their volunteer reviewers. ... Thus reviewers can
expect ACM to}

\begin{itemize}
\item {\em Maintain their anonymity}

  This has been \tods\ policy from the beginning.

\item {\em Ask them if they are willing to review before the submission is
  sent to them. The deadline for the review will accompany its request.}

  Reviewers are always asked first if they would be willing to review a paper.

\item {\em Provide guidelines on what constitutes a reviewing conflict of
  interest}

  There is a stated conflict of interest policy on the \tods\ web site (see
  {\tt http://www.acm.org/tods/}\-{\tt COI.html}). This policy is consistent with
  general ACM guidelines.

\item {\em Request them to review only submissions for which the editor
  feels they have expertise, and request only a limited number of reviews
  over the course of a year}

   The \tods\ Editorial Board very carefully chooses reviewers for submitted
   manuscripts.

   \tods\ strives to not overload referees. Specifically, \tods\ now has an
   explicit policy that referees will be expected to review at most one
   \tods\ paper in any twelve-month period.

   There are some provisos and exceptions to this policies. Informal
reviews and reviews of revised manuscripts can be quicker than two
months. Revised papers should be reviewed by the same referees, and this
review will probably occur within the twelve months, but that will just
extend the required interval before the next review. And referees are
welcome to volunteer for more reviewing than the maximum of one formal
review per year, if they wish.

\item {\em Recognize that they have the right to decline a requested review}

  The \tods\ Editorial Board policy is explicit on this right.

\item {\em Give a reasonable length of time for a review, where the
  particular length of time depends on the publication}

   \tods\ now has an explicit policy to allow at least two months for an initial
   formal review.

\item {\em Not routinely ask them to make up for delays introduced by other
  participants in the reviewing cycle}

  \tods\ has procedures in place to ensure that delays do not occur in the
  first place.

\item {\em Not ask them to provide reviews for submissions that do not
  satisfy either stated publications requirements (e.g., page count
  restrictions) or which are obviously inappropriate for the publication}

  The Editorial Board routinely desk rejects submissions that are
  inappropriate or that do not satisfy stated requirements.

\item {\em Acknowledge their efforts in the publication process, while
  maintaining confidentiality of which submissions they reviewed}

  \tods\ periodically publishes the collected names of reviewers.
  \tods\ employs single-blind reviewing. The identity of the reviewers of a
  manuscript will never be revealed to the authors or to the other reviewers.

\item {\em Inform them of the editorial decisions for the submission,
  including the author-visible portion of reviews}

  It is explicit \tods\ policy that the author-visible portion of reviews as
  well as the final editorial decision be provided to the reviewers once an
  editorial decision has been made. 

\item {\em Tell them who will see their review}

  Other than those just listed, no one else will be shown the reviews. 

\item {\em Recognize that reviewers own the copyright for their reviews.}

  This holds for all ACM publications.
\end{itemize}

\section{Libraries}

{\em In addition to reader rights, libraries can expect ACM to}

\begin{itemize}
\item {\em Provide institutional access to electronic versions at a
  reasonable price and take into account the needs of economically emerging
  and economically undeveloped countries}

  ACM has very reasonable institutional subscription rates, both for
  individual journals and for the ACM Digital Library as a whole. These
  rates are on a sliding scale.

\item {\em Enable fast access to the electronic versions throughout the
  world}

  The effort that has gone into providing fast access to individual
  subscribers also benefits patrons of libraries that have electronic subscriptions.

\item {\em Provide ongoing access, upon request, to electronic content to
  which a library has electronically subscribed, for the subscription
  period, should the subscription ever be canceled or should ACM remove
  titles from their electronic products, possibly for a fee}

  This provides assurance that libraries will always have access to the
  electronic content that they subscribed to.

\item {\em Respect fair use provisions of US copyright law}

  This has always been ACM policy.

\item {\em Allow libraries to fill interlibrary loan requests, for teaching,
  research, and other not-for-profit uses.}

  This has always been ACM policy.
\end{itemize}

\vspace*{-1ex}
\bibliographystyle{alpha}
\begin{thebibliography}{B82}

\bibitem[S02]{SN02}
Richard T.~Snodgrass, ``Progress on ACM's Becoming the Preferred Publisher,''
{\em Communications of the ACM}, 45(2):97--98, February, 2002.
{\tt http://doi.acm.org/10.1145/503124.503152}
\end{thebibliography}

\end{document}
